% Partie : histoire
% Chapitre : biographies
% Section : Gerhard Gentzen

\documentclass[../../../include/open-logic-section]{subfiles}

\begin{document}

\olfileid[fr]{his}{bio}{gen}

\olsection{Gerhard Gentzen}

\olphoto{gentzen-gerhard}{Gerhard Gentzen}

Gerhard Gentzen
est surtout connu
comme le fondateur
de la théorie
structurale de
la démonstration.
Il a introduit deux systèmes de !!{derivation},
la déduction naturelle et le calcul des séquents.
Il naquit le
24 novembre 1909
à Greifswald,
en Allemagne.
Il reçut
l'instruction
à domicile
pendant trois ans
avant d'entrer
dans une école
préparatoire,
où il avait
du retard
sur la plupart
de ses camarades.
Il se révéla
néanmoins un
élève brillant,
particulièrement
doué pour les
mathématiques.
Ses intérêts
étaient variés :
il écrivait
notamment des
poèmes pour
sa mère et
des pièces
pour le théâtre
de son école.

Gentzen commença
ses études
universitaires
à Greifswald,
puis étudia
à G\"{o}ttingen,
Munich et Berlin.
Il obtint
son doctorat
à l'université
de G\"{o}ttingen
en 1933,
sous la direction
de Hermann Weyl.
Paul Bernays
avait dirigé
l'essentiel
de son travail,
mais les nazis
l'avaient chassé
de l'université.
En 1935,
Gentzen devint
l'assistant
de David Hilbert.
Il exposa ces deux systèmes de !!{derivation}
dans ses articles
\emph{Untersuchungen
\"{u}ber das logische
Schlie\ss en I--II
[Recherches sur la
déduction logique
I--II]}.
En 1936,
il démontra
la cohérence
des axiomes
de Peano.

Les rapports
de Gentzen
avec les nazis
sont complexes.
Alors même
que son mentor
Bernays était
contraint de
quitter l'Allemagne,
Gentzen adhéra
à la section
universitaire
de la SA,
organisation
paramilitaire
nazie. Comme
beaucoup
d'Allemands,
il fut membre
du parti nazi.
Pendant la guerre,
il servit
comme officier
des télécommunications
dans une unité
de renseignement
aérien. En 1942,
il fut toutefois
libéré de ses
fonctions
à la suite
d'un effondrement
nerveux. On
ne sait pas
avec certitude
si ses convictions
le liaient
au parti nazi
ou s'il y
adhéra pour
assurer sa
réussite
universitaire.

En 1943,
Gentzen accepta
un poste
à l'Institut
de mathématiques
de l'Université
allemande
de Prague.
En 1945,
les habitants
de Prague
se soulevèrent
contre l'occupation
allemande.
Gentzen fut arrêté à Prague le 7 mai 1945 par des partisans.
Détenu à la prison de Karlovo náměstí, il fut astreint à des
travaux forcés.\footnote{Voir le témoignage de Franz Krammer,
transcrit par \href{https://dml.cz/dmlcz/139113}{Přemysl Vihan
(1993, p.~294)}.}
Il mourut
de malnutrition
dans sa cellule
le 4 août 1945,
à l'âge
de 35 ans.

\begin{reading}
Pour une biographie
complète de Gentzen,
voir
\citet{Menzler-Trott2007}.
\citet{Segal2014}
offre une lecture
intéressante sur
les mathématiciens
sous le régime
nazi, avec
un bref passage
sur la vie
de Gentzen.
Ses articles
sur la déduction
logique sont
disponibles
dans l'original
allemand
\citep{Gentzen1935a,Gentzen1935b}.
Leurs traductions
anglaises ont été
réunies dans un
volume
\citet{Gentzen1969}
qui comprend
aussi une
esquisse
biographique.
\end{reading}

\end{document}
